(i) is Theorem \claimref{dark-matter.mirror-tower-copy-one}(iii), with the equal asymmetry per mode of its clause (v) (the mirror instance of the former Premise P4), with Corollary \claimref{dark-matter.which-towers-can-carry-arrow} removing the odd towers from the asymmetric sum, read as a density ratio through the common dilution of Lemma \claimref{dark-matter.projection-every-towers-content-dilutes-ones}, and the half percent is Proposition \claimref{dark-matter.mirror-share-one-how-fragile}(i); $5.364 - 1 = 4.364$ and $4.364 \times 0.02237 = 0.0976$ are script-checked, with the survey uncertainties propagated as rational brackets. (ii) is the definition of a symmetric sector. (iii) is a scope statement resting on Corollary \claimref{dark-matter.which-towers-can-carry-arrow} and Proposition \claimref{dark-matter.no-coupling-rule-odd-towers-follows}; the rule it points to is used nowhere in this section and it asserts nothing about the remainder's value. The profile is Theorem \claimref{dark-matter.mirror-tower-copy-one}(i) for the mirror and Proposition \claimref{dark-matter.superselected-hence-dark-exactly-stable} for the coupling; the two stances follow from exact superselection and are costed in Section \ref{sec:premises}.

The other premises enter through the results cited: \claimref{dark-matter.p1-single-frozen-point} supplies the one metric, read as the sourcing rule of the common dilution; \claimref{dark-matter.p2-indifference} the equal asymmetry along the reversal orbit, the stored charge being a sector label; \claimref{synthesis.common-ceiling-coupling-premise} the common ceiling coupling along that orbit, on which the mirror's mass factor, and with it the share of one, rests; and \claimref{dark-matter.p3-nearest-dimension-representatives} the representatives at which Corollary \claimref{dark-matter.which-towers-can-carry-arrow} sorts the oriented towers.
